\subsection*{3.}
Aus dem eben Bewiesenen lassen sich Folgerungen für rationale Darstellung ziehen. Jede rationale absolute Invariante läßt sich --- wie durch Erweiterung von Zähler und Nenner mit den in bezug auf die Gruppe Konjugierten des Nenners unmittelbar einzusehen --- darstellen als Quotient von zwei, nicht notwendig teilerfremden, ganzen rationalen absoluten Invarianten. Daraus folgt, daß jede rationale absolute Invariante sich rational durch die Koeffizienten $G_{\alpha\alpha_1\ldots\alpha_n}(x)$ der Galoisschen Resolvente $\Phi(z,u)$ ausdrücken läßt. Dieser Satz läßt sich auch ohne Durchgang durch den Endlichkeitssatz auf verschiedene Art leicht beweisen; er findet sich schon formuliert in Weber II, \S\ 58.\NoetherSrcNote{*)}{Der dort gegebene Beweis ist allerdings nicht stichhaltig; es ist nämlich nur gezeigt, daß die Funktion $\Psi(t)$ in Formel (7) Invarianten zu Koeffizienten hat, nicht aber daß diese Invarianten rational durch die Koeffizienten von $\Phi(t)$ ausdrückbar sind. Diese Lücke wird vermieden, wenn man zur Darstellung der $x_i^{(k)}$ durch $\xi_k$ statt der Lagrangeschen Interpolationsformel ein bekanntes Differentiationsverfahren anwendet. Es ist dies die durch Differentiation der Identität $\Phi(-\xi_k,u)=0$ nach allen $u$ gewonnene Relation
\[
 \left[\frac{\partial\Phi}{\partial u_i}
      -x_i^{(k)}\frac{\partial\Phi}{\partial z}\right]_{z=-\xi_k}=0.
\]
Danach hat an Stelle von Formel (7) und (8) bei Weber für jede rationale Funktion $\omega(x)$ die folgende Formel zu treten:
\[
 \omega(x_1^{(k)},\ldots,x_n^{(k)})=
 \left.
 \omega\!\left(
   \frac{\partial\Phi/\partial u_1}{\partial\Phi/\partial z},\ldots,
   \frac{\partial\Phi/\partial u_n}{\partial\Phi/\partial z}
 \right)\right|_{z=-\xi_k},
\]
die nur noch Koeffizienten von $\Phi(z,u)$ enthält, und deren Summation über alle $k$ für Invarianten $\omega(x)$ die gewünschte Darstellung gibt.}
